\begin{center}
    \textsc{\huge On Credit Card Rewards}\\[0.1cm]
    \textsc{Monetizing Credit Rewards at Scale}\\[1em]
\end{center}

\section{PURPOSE}
\textit{We are building a credit card optimization platform for churners, travelers, and high-credit-score users by consolidating scattered knowledge into a graph of user-verified strategies.}

\textit{Our platform will help users plan credit card applications, track bonuses, and model outcomes--from their first card to their first free trip.}

\vspace{0.5em}
\section{PROBLEM AND OPPORTUNITY}
Credit card rewards are structured around the assumption that most users are inattentive. Banks offer large sign-up bonuses, knowing most inefficiently redeem them or carry interest-bearing balances. A small but growing segment of consumers approaches this system strategically: they apply for new cards regularly, plan spending to meet bonuses, and avoid interest entirely. This demographic relies on scattered tools, including:
\vspace{-0.5em}
\begin{itemize}[noitemsep]
    \item Spreadsheets to track sign-up bonuses
    \item Reddit and Discord for crowdsourced advice
    \item Manual estimates for sequencing and eligibility
\end{itemize}
\vspace{-0.5em}
While tools exist, none are purpose-built for the full lifecycle of a credit optimizer, especially one managing multi-card portfolios, facing denials, or mitigating shutdown risks from aggressive optimization.

\vspace{0.5em}
\section{SOLUTION AND EDGE}
We are building a platform that models credit card optimization as a sequence of decisions. Our platform:
\vspace{-0.5em}
\begin{itemize}[noitemsep]
    \item Pulls credit report and transaction data
    \item Tracks active cards, bonus progress, and eligibility for new cards
    \item Simulates reward paths and value accumulation over time
    \item Recommends next-best cards based on user-sourced strategy data
\end{itemize}
\vspace{-0.5em}
At our platform's core is a directional graph. Each node is a credit card (or relevant qualifier), and each edge is a user-verified path--for instance, going from a Discover Secured Credit Card to Amex Gold, with time between applications and spend history accounted for. Our directional graph grows increasingly robust as users interact with our platform and report real outcomes. 

\textit{No} existing card-monitoring tool embeds a community directly into their product, limiting organic growth and leaving untapped the potential of a social network where users share credit strategies.

\textbf{Competition:}
\vspace{-0.5em}
\begin{itemize}[noitemsep]
\item \textbf{MaxRewards} focuses on maximizing the expected value of cards a user \textit{already} possesses.\footnote{MaxRewards poses privacy and legal risks: it stores user banking credentials due to the lack of public, user-facing APIs (e.g., Amex, Chase), which triggers fraud alerts and account lockouts. It also automates all issuer rewards, including those not covered by the credit card, raising compliance concerns.}
\item \textbf{CardPointers} offers static\footnote{Emmanuel Crouvisier, the solo founder of CardPointers, manages the platform by his lonesome, including \textit{manually} updating card suggestions. CardPointers, however, boasts over 20,000 paying customers, and thus constitutes a ripe partnership opportunity: our architecture significantly outperforms CardPointers' in scalability and automation while Crouvisier has a six-year distribution moat.} suggestions that help users \textit{superficially} assess card eligibility and better apply spend to bonuses on opened cards.
\item Our platform integrates eligibility, sequencing, and long-term value modeling through a dynamic graph based on user-contributed outcomes.
\end{itemize}
\vspace{-0.5em}
Our \textit{ultimate} intent is to help users plan the timing and sequence of multiple credit card applications to maximize rewards and minimize risks over time.

\newpage
\section{MARKET SIZE AND STRATEGY}
\textbf{Market Estimate:}
\vspace{-0.5em}
\begin{itemize}[noitemsep]
    \item \textbf{TAM}: \$100M+, driven by consumer subscriptions and referral monetization
    \begin{itemize}
        \item 170M credit card users in the U.S.
        \item 48\% with FICO scores exceeding 750
        \item 10M with active interest in travel rewards, inferred from traffic to subreddits, blogs, and loyalty groups
        \item $\sim$1.75\% of transaction volume directed toward rewards (~\$100B USD)
    \end{itemize}
    
    \item \textbf{SAM:} \$10M--\$30M, targeting high-FICO users actively optimizing credit card rewards
    \begin{itemize}
        \item $\sim$5M users highly engaged with reward optimization tools and content
        \item Monetizable via \$2--\$6/month subscription or \$20--\$50/year referral value
    \end{itemize}
    
    \item \textbf{SOM:} \$1M--\$5M in near-term obtainable revenue
    \begin{itemize}
        \item Targeting $\sim$10K--25K early adopters through niche communities and influencer partnerships
        \item Initial revenue via premium subscription tiers and high-intent referrals
    \end{itemize}
\end{itemize}
\vspace{-0.5em}

\textbf{Initial user segments:}
\vspace{-0.5em}
\begin{itemize}[noitemsep]
\item Users already involved in churning via Reddit, Discord, and niche forums
\item Early-career professionals building credit and aiming to travel with points
\item Students with high credit scores and interest in travel or cashback rewards
\end{itemize}
\vspace{-0.5em}

\textbf{Growth plan:}
\vspace{-0.5em}
\begin{itemize}[noitemsep]
\item Direct outreach in high-engagement online forums
\item Socially engineer a contributor-driven community (e.g., rewarding influencers)
\item Gamified referral system based on user-verified card paths
\item Educational content focused on risk management, long-term planning, and successful strategies
\end{itemize}

\section{MONETIZATION AND OPERATIONS}
Our platform will offer a free trial, granting users full access to sophisticated planning tools, simulations, and data syncing for an introductory period. Post trial, premium features will remain available for \$4.99/month via subscription. Throughout, users can trust that all recommendations are based on objective analysis and never influenced by issuer incentives.

\textbf{Revenue streams include:}
\vspace{-0.5em}
\begin{itemize}[noitemsep]
\item \textbf{Subscription fees:} The primary source of revenue, offering paying users enhanced functionality and a seamless experience.
\item \textbf{Issuer referrals:} Used only when a referral aligns with the best, unbiased recommendation for the user—not as a default or intrusive strategy.
\item \textbf{Merchant partnerships:} Collaborations with merchants relevant to manufactured spending (e.g., prepaid and gift card providers), to create mutually beneficial referral opportunities.
\item \textbf{Graph data licensing:} Carefully anonymized churn path data may be licensed to fintech companies and academic researchers.
\end{itemize}

As our platform scales, additional roles in customer support, compliance, and partnership development will be added to support operations. 

Initial costs will center on critical infrastructure and early-stage growth:
\vspace{-0.5em}
\begin{itemize}[noitemsep]
\item Community data scraping (e.g., from public forums)
\item Integrations with essential data providers, such as Plaid and Experian
\item Cloud infrastructure to run simulations and deliver real-time services
\item Legal and compliance reviews to ensure regulatory alignment
\item Budget for community development, including moderation resources and user rewards
\end{itemize}
\vspace{-0.5em}

In conclusion, this structure lays the groundwork for a fundamentally new way users can unlock the full value of their credit.
